\section{从全局状态到局部行为}

\begin{proposition}[嵌套状态解释]
如果全局意识状态依赖于有结构的动力学组织，那么局部行为框架就可以被解释为更广泛生物状态景观中的、任务特定的嵌套机制。
\end{proposition}

\begin{discussion}
这一命题是刻意保守的。
强证据\emph{并不}意味着阅读一条定理与进入睡眠是同一个事件。
它真正支持的是一套共享词汇：稳定性、转变、协调与机制变化。
随后，本书再进行一个\textbf{中等}强度的解释性推进，把局部行为框架视为这一更广泛景观中的小尺度机制。
这一推进之所以有用，是因为它让姿势、注意力、环境与提示线索可以被讨论为机制管理变量，
而不再只是模糊的情绪描述。
\end{discussion}